\chapter{Capstone: from requirement to operated release}

\section{The assignment}
Design and implement a service that imports a customer order file, validates it,
stores accepted orders, exposes a query endpoint, and produces an operational
report. You may use the repository's Ledgerline examples as a starting point,
but the capstone should make its own decisions explicit.

\section{Required deliverables}
\begin{enumerate}
  \item a one-page problem statement with stakeholders, non-goals, and risks;
  \item an API and data specification with invalid examples;
  \item an architecture diagram and at least two ADRs;
  \item a working implementation with unit, component, and migration tests;
  \item a threat model and security verification checklist;
  \item a container image and CI workflow;
  \item an SLO, dashboard or query plan, and recovery runbook;
  \item a migration plan for one schema change;
  \item an incident scenario and blameless postmortem;
  \item a short reflection on accessibility, internationalization, sustainability,
        and AI-assisted development.
\end{enumerate}

\section{A release checklist}
\noindent\begin{longtable}{@{}p{0.07\textwidth}p{0.47\textwidth}p{0.28\textwidth}@{}}
\toprule
Done & Check & Evidence \\
\midrule
 & Behaviour and non-goals are clear & reviewed requirements \\
 & Invalid input is bounded and rejected & tests and validation code \\
 & Authorization is tested on allowed and denied paths & security tests \\
 & Data invariants exist at the right boundary & schema and migration checks \\
 & Retries and unknown outcomes are specified & contract and state model \\
 & Artifact is traceable to source and checks & CI record \\
 & Deployment is observable and reversible & runbook and rollout plan \\
 & Backups or exports can be restored & restore evidence \\
 & Accessibility and locale paths are tested & manual or automated evidence \\
 & Dependencies and licenses are reviewed & inventory and policy \\
\bottomrule
\end{longtable}

\section{How to present the design}
Begin with the user-visible outcome. Show one complete path, then one failure
path. State assumptions before diagrams. Explain why the design is sufficient
for the current constraints and what would force a different architecture.
Demonstrate a test failure or simulated incident rather than claiming that a
feature is “production ready” without evidence.

\section{Final perspective}
The central habit of software engineering is to turn invisible assumptions into
visible, testable, revisable decisions. Reliable systems are not produced by one
perfect design. They emerge from boundaries that fit the problem, feedback that
arrives before consequences become large, and teams that treat operation and
maintenance as part of building.

\begin{keyidea}
Build the smallest system that satisfies the important constraints, but make its
constraints explicit enough that a future team can safely discover when “smallest”
is no longer sufficient.
\end{keyidea}
